% Figure 5. Gray fill is the side phi < 0; the boundary is tangent at 0.
\begin{tikzpicture}[field/diagram, x=1.15cm, y=1.0cm]
    \def\LocalBoundary{
        (-1.75,1.12)
        .. controls (-1.05,0.76) and (-0.65,0) .. (0,0)
        .. controls (0.65,0) and (1.1,0.72) .. (1.75,1.35)
    }
    \fill[FieldRegion] \LocalBoundary
        -- (1.75,1.85) -- (-1.75,1.85) -- cycle;
    \draw[field/secondary] (-1.75,-1.50) rectangle (1.75,1.85);
    \draw[field/axis] (-2.15,0) -- (2.45,0)
        node[right] {$x_n=0$};
    \draw[field/axis] (0,-1.7) -- (0,2.2) node[above] {$x_n$};
    \draw[field/curve] \LocalBoundary;
    \node at (-0.82,1.40) {$\phi<0$};
    \node at (-0.85,-0.77) {$\phi>0$};
    \node[below right, inner sep=3pt] at (0,0) {$(0,0)$};
    \node[right] at (2.0,1.15) {$\mathbb R^{n-1}\times\mathbb R$};
    \draw[field/guide] (-1.89,-1.5) -- (-2.24,-1.5)
                       (-1.89,1.85) -- (-2.24,1.85);
    \draw[field/dimension] (-2.18,-1.5) -- (-2.18,1.85)
        node[midway, field/label] {$V$};
    \draw[field/guide] (-1.75,-1.62) -- (-1.75,-2.02)
                       (1.75,-1.62) -- (1.75,-2.02);
    \draw[field/dimension] (-1.75,-1.94) -- (1.75,-1.94)
        node[midway, field/label] {$U$};
\end{tikzpicture}
